% 这是示例论文骨架，不是默认 TeX 入口。
% 默认排版模板是 templates/tex/main.tex；正文章节应由题目和证据决定。
% 不要把下列标题机械填进每一篇论文。

\section{问题重述与任务分解}
\subsection{问题背景}
% 只保留理解后续模型所需的背景、对象关系和实际约束。

\subsection{问题重述}
% 将实际子问题改写为可检验的任务，明确输入、输出、目标、约束和评价指标。

\section{问题分析}
% 用连续段落说明对象关系、关键张力和各问之间的依赖；不要用清单代替论证。

\section{数据、假设与符号}
\subsection{数据来源与预处理}
\subsection{模型假设}
% 假设、算法步骤等天然并列的对象可以使用列表。
\subsection{符号说明}

\section{子问题的模型与求解}
% 按实际题目改写标题；不存在的子问题不要保留。

\section{结果验证与稳健性}

\section{讨论、局限与结论}
% 结论用连续段落回答题目，不要把路线复述再列一遍。
